\chapter{Data Dependencies, Scratch Directories, and Generated Launchers}
\label{ch:lessons-03-04}

\begin{lessonbox}{Lecture 08 --- Data dependencies with \texttt{workflow://}}
\textbf{Objective.} Let DAGonStar infer dependencies from data references.

\textbf{Concept.} \texttt{workflow:///task/path} references a file inside
another task's scratch directory. \texttt{Workflow.make\_dependencies()}
parses this reference and creates the dependency.

\textbf{Course example} (\texttt{examples/08-dataflow/workflow\_references.py}, Lecture~08).
\lstinputlisting[caption={Lecture 08: \texttt{workflow\_references.py}},label={lst:lecture-08}]{../../examples/08-dataflow/workflow_references.py}

\textbf{Run} from the repository root:
\begin{lstlisting}[style=bash]
python examples/08-dataflow/workflow_references.py
\end{lstlisting}

\textbf{Recorded output} (\texttt{examples/08-dataflow/output/workflow\_references.txt}):
\lstinputlisting[style=bash]{../../examples/08-dataflow/output/workflow_references.txt}

\textbf{Scientific practice note.} Dataflow references make scientific
provenance easier to read: an edge exists because a named data product is
consumed.
\end{lessonbox}

\section{Discussion: Inference versus Assertion}

Lecture~08 is the data-oriented twin of Lecture~07, and the contrast is
instructive when read side by side. In Lecture~07, a reader of the source code
must already know, extrinsically, that ``validate parameters'' precedes
``launch the model.'' In Lecture~08, the dependency is \emph{visible inside the
command string itself} --- \texttt{workflow:///observations/data/series.csv}
simultaneously documents \emph{and} enforces the relationship. This is
precisely the property the documentation calls easier-to-read scientific
provenance: the graph edge exists \emph{because} a named data product is
consumed, not by separate assertion.

\begin{lessonbox}{Lecture 10 --- Scratch, launchers and logs}
\textbf{Objective.} Understand what DAGonStar writes during execution.

\textbf{Concept.} Each task has a working directory and an internal
\texttt{.dagon} directory containing generated scripts and captured output.

\textbf{Course example} (\texttt{examples/10-scratch/inspect\_scratch.py}, Lecture~10).
\lstinputlisting[caption={Lecture 10: \texttt{inspect\_scratch.py}},label={lst:lecture-10}]{../../examples/10-scratch/inspect_scratch.py}

\textbf{Run} from the repository root:
\begin{lstlisting}[style=bash]
python examples/10-scratch/inspect_scratch.py
\end{lstlisting}

\textbf{Recorded output} (\texttt{examples/10-scratch/output/inspect\_scratch.txt}):
\lstinputlisting[style=bash]{../../examples/10-scratch/output/inspect_scratch.txt}

\textbf{Scientific practice note.} Generated launchers are a practical
provenance artifact. They show the actual command DAGonStar executed after
preprocessing.
\end{lessonbox}

\section{Anatomy of a Task Working Directory}

\Cref{fig:task-dir} formalizes the layout that Lecture~10 asks the reader to
discover by inspection. This layout recurs, with backend-specific additions,
across every execution backend treated in Part~II.

\begin{figure}[h]
\centering
\begin{tikzpicture}[
  dir/.style={rectangle, draw=dagonblue, thick, fill=blue!4,
    minimum width=58mm, minimum height=8mm, font=\ttfamily\small, anchor=west},
  sub/.style={rectangle, draw=dagonaccent, thick, fill=white,
    minimum width=50mm, minimum height=7mm, font=\ttfamily\footnotesize, anchor=west}
]
  \node[dir] (root) at (0,0) {/tmp/\textit{\textless timestamp\textgreater}-inspect/};
  \node[sub] (out) at (0.6,-1.1) {product.txt \quad \footnotesize(task output)};
  \node[sub] (dagon) at (0.6,-2.0) {.dagon/};
  \node[sub] (launcher) at (1.2,-2.9) {launcher.sh \quad \footnotesize(generated, executed command)};
  \node[sub] (context) at (1.2,-3.7) {context.sh \quad \footnotesize(execution context/environment)};
  \node[sub] (stdout) at (1.2,-4.5) {stdout.txt \quad \footnotesize(captured standard output)};
  \node[sub] (inputs) at (1.2,-5.3) {inputs/\textit{...} \quad \footnotesize(staged \texttt{workflow://} files, if any)};
\end{tikzpicture}
\caption{Canonical layout of a task's scratch (working) directory, as
discoverable through Lecture~10's verification commands.}
\label{fig:task-dir}
\end{figure}

\begin{practicebox}
Treat \texttt{.dagon/launcher.sh} as the closest thing the engine offers to a
laboratory notebook entry for a single task: it is the literal, fully
resolved command that ran, \emph{after} every \texttt{workflow://} reference
has been rewritten to a concrete staged path. When debugging a failed
scientific task, read this file before re-reading the Python source --- it
removes one layer of indirection between ``what I wrote'' and ``what
actually executed.''
\end{practicebox}

\section{From Inspection to Debugging Practice}

Lecture~10's verification procedure --- listing the task directory, then
reading \texttt{launcher.sh}, \texttt{context.sh}, and \texttt{stdout.txt}
in sequence --- is presented in the tutorial as a one-time exercise in
familiarization. It is worth elevating it explicitly to a general-purpose
debugging procedure, because the same three files, read in the same order,
are the correct first response to \emph{any} unexpected task failure
encountered later in this book, including the heterogeneous backends of
Part~II.

\begin{enumerate}[leftmargin=1.8em]
  \item \textbf{Read \texttt{launcher.sh} first.} This file contains the
    fully resolved command, after every \texttt{workflow://} reference has
    been rewritten to a concrete, staged path (\Cref{ch:schema}). The most
    common class of early failure --- a referenced file that does not
    exist, or a path that resolved to somewhere unexpected --- is usually
    visible immediately at this step, before any need to inspect output at
    all.
  \item \textbf{Read \texttt{context.sh} second.} This file captures the
    execution context, including environment-related setup the launcher
    relies on; discrepancies between the environment a task actually saw
    and the environment a developer assumed are a frequent source of
    ``works on my machine'' failures, particularly once
    \Cref{ch:external-software}'s environment-modules and virtual-environment
    activation patterns are involved.
  \item \textbf{Read \texttt{stdout.txt} last.} Only once the command and
    its execution context are confirmed correct does the captured standard
    output --- which records what the external program itself reported ---
    become the most informative artifact.
\end{enumerate}

\begin{practicebox}
This ordered procedure exists because each file answers a strictly more
specific question than the one before it: \texttt{launcher.sh} answers
``what command, exactly, did DAGonStar decide to run?''; \texttt{context.sh}
answers ``in what environment did it decide to run that command?''; and
\texttt{stdout.txt} answers ``what did the command itself report once it
ran?'' Skipping directly to \texttt{stdout.txt} --- the most tempting file to
open first, since it most resembles ordinary terminal output --- frequently
wastes time chasing a symptom (a program's own error message) whose root
cause (a wrong path, a missing dependency, an unset environment variable)
was already visible one file earlier.
\end{practicebox}

\section{The Provenance Value of an Immutable Launcher}

A subtlety worth making explicit, building on \Cref{ch:schema}'s discussion
of command rewriting: \texttt{launcher.sh} is generated \emph{once}, before
the task executes, and is not regenerated or modified afterward. This
immutability is what makes it trustworthy as a provenance artifact in the
sense already developed in \Cref{ch:schema}'s discussion of the provenance
graph. A researcher inspecting a scratch directory days or months after a
workflow run can be confident that \texttt{launcher.sh} reflects exactly the
command that executed at the time of the run --- not a reconstruction, not
an approximation, and not a file that some later process might have
silently altered.

Contrast this with the alternative of relying solely on the original Python
workflow source to reconstruct what ran. The Python source records the
\emph{template} a command was built from --- including any unresolved
\texttt{workflow://} references --- but not the runtime values that
template was eventually instantiated with, nor whether the workflow was
later modified and rerun with different parameters. Two facts can only be
recovered from \texttt{launcher.sh} itself: which specific staged path each
reference resolved to at execution time, and the precise, final form of any
command-line argument that depended on a runtime value. This is the concrete
sense in which the generated launcher, not the source script, is the
authoritative record of what a given scratch directory actually represents.

\begin{exercise}
Run Lecture~08's workflow twice in succession, without deleting the resulting
scratch directories between runs. Compare the two generated
\texttt{statistics} task directories' \texttt{launcher.sh} files. What
differs between the two, and what does this tell you about which parts of a
generated launcher are stable across repeated runs of an unmodified
workflow, versus which parts (if any) vary?
\end{exercise}

\section{Chapter Summary}

Lectures~08 and~10 together complete the data-oriented half of dependency
modeling and expose the concrete, on-disk artifacts the engine produces while
enacting it: a staged input file under \texttt{.dagon/inputs/}, a rewritten
\texttt{launcher.sh}, and a captured \texttt{stdout.txt}. This chapter
additionally elevated the three-file inspection procedure of Lecture~10 into a
general debugging discipline applicable throughout the rest of this book,
and argued for the generated launcher's immutability as the structural
reason it functions as an authoritative provenance record, independent of
and complementary to the original workflow source code. \Cref{ch:lessons-05-07}
turns to the failure and recovery mechanisms layered on top of this same
on-disk model: cycle detection, checkpointing, and configurable data staging.
